% ============================================================================
\section{Discussion}\label{sec:discussion}
% ============================================================================

\subsection{Implications for guidance design}
\begin{enumerate}[leftmargin=*,itemsep=3pt]
\item \emph{Do not expect a single LCvx-type SOCP for attitude-rate-limited landing.} Within the class of convex programs on the
physical state with affine data entry, exactness is impossible, and this cannot be fixed by adding slacks, cuts, lifted
variables, discretization or a line search over $t_f$ (\cref{thm:A,thm:Bstar,lem:LS,thm:D}); the state-constrained core of this statement is machine-checked. Any such program, if it contains all
genuine trajectories, accepts boundary data that no flyable trajectory can land from. The flight software then has no way to
detect infeasibility from the convex program alone. The statement is about exactness. It does not say that a single SOCP cannot be
\emph{sound}; see item~3.
\item \emph{Outer attitude loop, exact convex inner problem.} The nonconvexity lies entirely in the choice of attitude profile
(\cref{thm:CU}). For a fixed attitude profile, the time-varying-cone LCvx is exact (\cref{thm:CW,thm:CS}). On a pure-TVC vehicle the
problem given the attitude is \emph{exactly} convex, with no relaxation (\cref{thm:C6}). A natural architecture is therefore an
outer search over a low-dimensional family of attitude profiles (or a sequential update of $b$) around an exact convex inner
problem. Every inner solution is then genuinely flyable, and the inner problem certifies infeasibility for the given profile.
\Cref{sec:practice} implements this architecture and evaluates it numerically. Standard successive convexification, which
linearizes the attitude together with everything else, performed as well, so the architecture's value is speed and flyable
intermediate iterates rather than better plans.
\item \emph{Or no loop at all: a sound single SOCP.} On a pure-TVC vehicle the outer loop exists only because the problem is posed
in centre-of-mass coordinates, where attitude and thrust multiply. In centre-of-percussion coordinates the dynamics are linear, and
restricting the actuator limits instead of relaxing them gives one SOCP whose plans are all flyable at every instant (\cref{thm:F}).
What it gives up is measured in \cref{sec:practiceF}: $0.1\%$ of fuel at the operating point and a few tenths of a second of
minimum flight time, concentrated in manoeuvres that tilt in order to shed thrust. Because it needs no initial trajectory, trust
region or iteration budget and solves in tens of milliseconds, it can replan at the control system's own time scale, and in our
simulations that mattered more for touchdown accuracy than the last fraction of a percent of optimality. For a new gimballed vehicle
we would start here and keep successive convexification for the manoeuvres this program refuses.
\item \emph{Count computation time.} Closed-loop comparisons that apply each plan at the instant its state was sampled favour slow
planners. With plans applied after their solver time, the planner built on \cref{thm:C6} lost 4 of 40 in-air landings it had made
before, and the lateral error of successive convexification roughly doubled.
\item \emph{Certificates.} The degree-2 SDP lift gives a rigorous lower bound on the minimum landing error and certifies
infeasibility of exact landing when positive. Its global optimality can be certified a~posteriori (rank 3, stress test), and
a~priori for coplanar targets (\cref{thm:stress,thm:coplanar}). It is a useful companion to a local or sequential solver, not a
replacement.
\item \emph{How much does it matter?} For in-air descents of a small lander, the exhibited nonconvexity is centimetre-scale, and
a sound convex program loses little in value ($0.1\%$ of fuel in \cref{sec:practiceF}). The large effects are elsewhere. First, near
lift-off (pad-launched hops) the landing set can have a metre-scale hole. Second, 3-DOF LCvx solutions ask for thrust-direction
changes far beyond the vehicle's angular acceleration (\cref{sec:numerics}). Third, on a pure-TVC vehicle the 3-DOF optimum is
essentially never exactly flyable. The ``upright body plus pointing cone'' 3-DOF model has effective cone zero.
\item \emph{Tilt-angle rate limits.} Constraining only the tilt angle from vertical, with free azimuth, removes rigidity and
appears to give convex landing sets numerically. It is not SOC-representable either, and it is not the physical constraint
of a rigid body.
\end{enumerate}

\subsection{Open problems}
\begin{enumerate}[leftmargin=*,itemsep=3pt]
\item \emph{Pure TVC, upright landing at rest, with altitude constraint.} This is the actual small-lander case. It is proved, computer-assisted and
certified in exact arithmetic, with constant mass and unlimited gimbal rate, at $t_f=1,2$~s without a tilt cap and at $t_f=3,5$~s in
the tilt-capped class (\cref{thm:TVC1}). The fully realistic model (rate-limited gimbal, mass depletion, free roll) is nonconvex
numerically at every horizon tried ($t_f=3$--$7$~s, gaps $0.5$--$3.3$~cm; a 1.7~m hole at hop data). Still open: certificates with a
gimbal-rate limit, with mass depletion, at $t_f=7$~s capped (fails by 4.8\%), for uncapped $t_f>2.3$~s (tumbling lemma), for the box
gimbal and free roll beyond $t_f=1$~s, and with free final time.
\item \emph{Rigor at realistic constants.} For our vehicle, Lemma S (hard soft landing, per instance) and the hop-instance hole
are numerical. The available quadratic, Wirtinger and S-procedure certificates are provably too loose. A two-dimensional
rigidity argument or validated branch-and-bound is needed. The $\rho_1=0$ landing-set result with altitude constraint is open at
$T/W\approx1.27$.
\item \emph{CSD2$'$.} Is the degree-2 lift exact in continuous time with fixed thrust, or for the convexified-lens and $\delta=0$
discretizations? Is there a tractable exact lift with variable throttle? (The natural degree-2 throttle lift and the separate-axis
lift CSD3 were both found lossy numerically.)
\item \emph{Minimum-fuel value-losslessness.} Does $V$ agree with its convex envelope $\conv V$ on $F$ for minimum fuel? Theorem A-opt reduces this to a convexity
question, and numerics show no violation. A positive answer would still not give an LCvx in the detection-sound sense.
\item \emph{Long horizons and coning.} At long horizons with tilt caps, coning fills most but not all of the dents. Whether
nonconvexity persists for all horizons in the tent variant, and at faces with long capped plateaus, is open.
\item \emph{Quantitative state-constrained gaps.} The rigorous B$^*$ gap ($\sim10^{-5}$~m) is far below the numerical one
($\sim10^{-2}$). A sharp quantitative version would settle how much a guidance designer actually loses.
\item \emph{The sound program of \cref{thm:F}.} (a) Formalize it; every lemma is elementary. (b) Bring the gimbal and throttle slew
limits inside the guarantee: bounds with worst-case constants are valid but too loose to be useful, and a restriction with
reference parameters, as in rows (R3)--(R5), is the natural route. (c) Bound the loss: \cref{prop:F-inc} says the program is
incomplete, and the experiments say by how much on one vehicle, but there is no theorem bounding the fuel or flight-time penalty.
(d) Remove \cref{ass:cop}: with $\varepsilon$ varying along the flight the output $y$ is no longer flat, and an extra term
$2\dot\varepsilon\dot b$ appears in $\ddot y$. (e) Is there a larger convex sound set, for instance one that allows tilting to shed
thrust in a fixed direction?
\end{enumerate}

\paragraph{Reproducibility.} \intshare{Proof notes (\texttt{math/}), the Lean project (\texttt{lean/}), numerical scripts and outputs
(\texttt{numerics/}), the planners and their benchmarks (\texttt{solver/}) and the claims registry (\texttt{comms/CLAIMS.md}) are kept
in the project repository. \texttt{paper/CLAIMS-MAP.md}
maps every theorem of this paper to its proof file, registry row and Lean name.}{The supplementary archive contains the Lean project
(\texttt{lean/}: sources, \texttt{Verify.lean}, and the notes \texttt{MODEL.md}, \texttt{POSITIVE.md}, \texttt{SLICE.md},
\texttt{SOLVER.md} and \texttt{TVC.md}, which map each Lean definition to its mathematical meaning) and the scripts, data and
derivation notes of the computer-assisted pure-TVC certificates (\texttt{certificates/}), and the code, tests and proof note of the
sound single convex program of \cref{sec:F} (\texttt{landing/}). Its \texttt{README} gives the commands that reproduce every check.}
